% !TeX root = main.tex
% Synchronized with 00_摘要.md v0.7; Tsai format.
\begin{abstractzh}
\setcounter{page}{1}

本研究旨在建立一套適用於動態環境之移動機械臂預測導航與受約束操作系統，整合環境感知、幾何模型預測控制與底盤—手臂協調機制，以處理障礙避讓、末端追蹤及操作可達性等問題，並探討機器人在持續移動下執行操作任務之控制方法。

導航部分以平面剛體運動群 \(SE(2)\) 描述底盤位姿與追蹤誤差，透過有限時域預測規劃控制輸入，並結合障礙運動預測、控制屏障條件及原始雷射掃描防護，處理動態環境中的路徑追蹤與避障。

協同操作部分針對位置命令介面建立手臂執行模型，以關節位置與設定點構成增廣狀態，描述命令輸入至實際關節運動的追蹤動態。控制設計依據底盤與手臂對末端運動的共同影響，於預測時域內聯合最佳化控制輸入，並納入末端追蹤目標、關節限位與輪級運動限制。系統另以視覺量測建立接近與操作參考，依任務階段銜接導航與協同操作。

本研究以動態導航與抽屜操作模擬，評估所設計方法之追蹤表現、任務完成情形與執行效率。

\keywordzh{\mbox{移動機械臂}、\allowbreak\mbox{幾何模型預測控制}、\allowbreak\mbox{全身協調控制}、\allowbreak\mbox{執行動態建模}、\allowbreak\mbox{控制屏障函數}、\allowbreak\mbox{移動中操作}}
\end{abstractzh}
